\input{../tex-inputs/latex-leseansicht-vorspann}

\section[Arthur Schnitzler an Gustav Schwarzkopf, {[}10. 10. 1905?{]}]{L04195 Arthur Schnitzler an Gustav Schwarzkopf, {[}10. 10. 1905?{]}}
\nopagebreak\mylabel{L04195v}
\rehead{ }\normalsize\beginnumbering\briefempfaengerindex{Schwarzkopf, Gustav@\textsc{Schwarzkopf, Gustav}!zzzSchnitzler, Arthur@\emph{von Arthur Schnitzler}!1905-10-102@{{[}10. 10. 1905?{]}}|(be}
\toendnotes[C]{\smallbreak\pagebreak[2]}
\correspDesc{Versand  durch Arthur Schnitzler am [10. 10. 1905?] in Wien
\newline{}Erhalt  durch Gustav Schwarzkopf im Zeitraum [11. 10. 1905 – 15. 10. 1905?] in Tiefer Graben 23}\toendnotes[C]{\smallbreak}
\Standort{CUL, Schnitzler, B 96.}
\physDesc{Briefkarte, 121 Zeichen
\newline{}Handschrift: Bleistift, deutsche Kurrent}\toendnotes[C]{\smallbreak}
\pstart
           \noindent{}{\pb}Lieber Guſtav. Hier iſt der \label{K_L04195-1v}\edtext{Sitz\eventindex{Burgtheater@\textbf{Burgtheater}!Uraufführung von Zwischenspiel, 12.10.1905@Uraufführung von Zwischenspiel, 12.10.1905|pwv}}{\lemma{\textnormal{\emph{Sitz}}}\Cendnote{\textnormal{Die Karte ist undatiert und es sind
               viele Tage möglich, an denen Schwarzkopf\pwindex{Schwarzkopf, Gustav 7.\,11.\,1853 Wien – 13.\,11.\,1939 ebd.@\textsc{Schwarzkopf, Gustav} (7.\,11.\,1853 Wien – 13.\,11.\,1939 ebd.)|pwk} nach dem 
                  Theater zu Louise Schnitzler\pwindex{Schnitzler, Louise 8.\,7.\,1840 Kőszeg – 9.\,9.\,1911 Wien@\textsc{Schnitzler, Louise} (8.\,7.\,1840 Kőszeg – 9.\,9.\,1911 Wien)|pwk} gekommen sein kann. Schränkt
                  man ein, dass es sich um einen Theaterbesuch nach einer Aufführung eines Stückes
                  von Schnitzler handelte, so schränken sich die
                  Möglichkeiten ein und es dürfte sich um die 
                  Uraufführung von Zwischenspiel\eventindex{Burgtheater@\textbf{Burgtheater}!Uraufführung von Zwischenspiel, 12.10.1905@Uraufführung von Zwischenspiel, 12.10.1905|pwk} am 12. 10. 1905
                  handeln, bei der in Folge in der Frankgasse 1\oindex{Wien@\textbf{Wien}!IX., Alsergrund@\textbf{IX., Alsergrund}!Frankgasse 1@\textbf{Frankgasse 1}|pwk} gefeiert wurde.
                  Das fügt sich auch ein in die Aussage des Briefes vom XXXX Auszeichnungsfehler: Dokument L04019 nicht gefunden,
                  wonach Schnitzler{ }Schwarzkopf\pwindex{Schwarzkopf, Gustav 7.\,11.\,1853 Wien – 13.\,11.\,1939 ebd.@\textsc{Schwarzkopf, Gustav} (7.\,11.\,1853 Wien – 13.\,11.\,1939 ebd.)|pwk}
                  nicht in seiner Loge unterbringen kann, aber auf seine Anwesenheit setzt. Die Karten dürfte
                  Schnitzler wiederum am Dienstag, dem 10. 10. 1905 erhalten haben,
                  vgl. XXXX Auszeichnungsfehler: Dokument L01559 nicht gefunden.}}}\label{K_L04195-1} und
      es bleibt dabei, dſs Sie da{\geminationn} bei uns
      eſſen. Meine Mama\pwindex{Schnitzler, Louise 8.\,7.\,1840 Kőszeg – 9.\,9.\,1911 Wien@\textsc{Schnitzler, Louise} (8.\,7.\,1840 Kőszeg – 9.\,9.\,1911 Wien)|pwv} bittet Sie
      darum.\pend
           \pstart Herzlich Ihr \spacefill\mbox{Arthur}\pend{}\selectlanguage{ngerman}\endnumbering\briefempfaengerindex{Schwarzkopf, Gustav@\textsc{Schwarzkopf, Gustav}!zzzSchnitzler, Arthur@\emph{von Arthur Schnitzler}!1905-10-102@{{[}10. 10. 1905?{]}}|)be}\mylabel{L04195h}
\begin{anhang}
\end{anhang}\newcommand{\dateiname}{L04195}\newcommand{\titel}{Arthur Schnitzler an Gustav Schwarzkopf, [10. 10. 1905?]}\newcommand{\editorInnen}{Herausgegeben von Jahnke, SelmaMüller, Martin Anton}\input{../tex-inputs/latex-leseansicht-abspann}
